\section{Worked convergence, cost and scale calculations}
\label{sec:numerical}

\subsection{An explicit favorable transfer model}

Use \eqref{eq:quadratic} with $H=\diag(1,10)$ and $\alpha=0.05$.
Let scratch error be $(1,1)$ and transferred error be $(1/2,1)$.
The two contraction factors are $0.95$ and $0.5$, so
\begin{align}
 f_0(n)-f_*&=\frac12\left(0.95^{2n}+10\,0.5^{2n}\right),\\
 f_G(n)-f_*&=\frac12\left(\frac14\,0.95^{2n}
                                  +10\,0.5^{2n}\right).
 \label{eq:toycurves}
\end{align}
The transfer reduced only the slow-mode initial error.
For target excess loss $\varepsilon=0.01$, evaluate the exact expressions:
\begin{center}
\begin{tabular}{@{}llrl@{}}
\toprule
Initialization & Update $n$ & Excess loss & Target reached?\\
\midrule
Scratch & 38 & 0.01013827358 & No\\
Scratch & 39 & 0.00914979190 & Yes\\
Transferred & 24 & 0.01065719879 & No\\
Transferred & 25 & 0.00961812191 & Yes\\
\bottomrule
\end{tabular}
\end{center}
Both squared factors lie strictly between zero and one, so loss is
monotonically decreasing. The neighboring checks therefore establish
the \emph{first} hitting updates:
\begin{equation}
 N_0=39,\qquad N_G=25,\qquad
 S_{\rm language}=\frac{39}{25}=1.56.
\end{equation}
These values come directly from powers in \eqref{eq:toycurves}, not
from a simulated optimizer.

In a single dominant mode, with factor $\rho=|1-\alpha\lambda|\in(0,1)$,
the exact mode loss is $A\rho^{2n}$. For $A>\varepsilon$,
\begin{equation}
 N(A,\varepsilon)=
 \left\lceil\frac{\log(\varepsilon/A)}{2\log\rho}\right\rceil.
 \label{eq:onemodehit}
\end{equation}
Reducing the coefficient amplitude by $r\in(0,1)$ replaces $A$ by
$r^2A$. Before integer rounding and any zero-step clipping, the saved
updates are
\begin{equation}
 N(A,\varepsilon)-N(r^2A,\varepsilon)
 \simeq\frac{\log r}{\log\rho}.
\end{equation}
For $r=1/2$ and $\rho=0.95$, this is about 13.5 updates, consistent with
the exact 14-update difference above after rounding and the fast mode.
The formula explains why reducing slow-mode error matters.

\subsection{Break-even with synthetic cost}

Against scratch, the necessary and sufficient cost inequality under
\eqref{eq:twocosts} is
\begin{equation}
 c_0^{\rm lang}N_0-c_G^{\rm lang}N_G
 >C_G^{\rm gen}+c_G^{\rm ppt}K.
 \label{eq:breakeven}
\end{equation}
If language-step costs are equal to $c$,
\begin{equation}
 N_0-N_G>
 \frac{C_G^{\rm gen}}c+\frac{c_G^{\rm ppt}}cK.
\end{equation}
For equal-cost steps and zero generation cost, the toy model gives
\begin{align}
 K=10:&\quad S_{\rm total}=\frac{39}{10+25}
                    =1.114285714,\\
 K=20:&\quad S_{\rm total}=\frac{39}{20+25}
                    =0.866666667.
\end{align}
The same $1.56\times$ language-stage speedup can therefore correspond to
either a total gain or a total loss. With $K=10$, generation must cost
less than four language-step units to retain a strict total benefit.
At exactly four units, total costs tie.

For graph versus grid, include the \emph{difference} in preparation
cost rather than automatically charging a common cost twice:
\begin{equation}
 c_{\rm grid}^{\rm lang}N_{\rm grid}
       -c_G^{\rm lang}N_G
 >
 (C_G^{\rm gen}-C_{\rm grid}^{\rm gen})
 +(c_G^{\rm ppt}-c_{\rm grid}^{\rm ppt})K.
\end{equation}
Caching changes amortized generation cost per use. It does not remove
the initial cost. A stated conservative full-cost metric and an amortized
metric can both be reported, provided their denominators are explicit.

\subsection{Exact binary-pilot transformer parameter count}

Let vocabulary size be $V$, width $D$, context length $W$, and block
count $L_b$. A block has attention weights $3D^2+D^2$, feed-forward
weights $4D^2+4D^2$, linear biases $3D+D+4D+D=9D$, and two layer
normalizations with total $4D$ scale/shift parameters. Thus
\begin{equation}
 P_{\rm block}=12D^2+13D.
\end{equation}
For tied token/output embeddings and a final layer normalization,
\begin{equation}
 P_{\rm model}=VD+WD+
             L_b(12D^2+13D)+2D.
 \label{eq:paramcount}
\end{equation}
Substituting $V=274,D=192,W=128,L_b=4$:
\begin{align}
 VD&=52\,608,\qquad WD=24\,576,\\
 P_{\rm block}&=12(192)^2+13(192)=444\,864,\\
 P_{\rm model}
 &=52\,608+24\,576+4(444\,864)+384\\
 &=\boxed{1\,857\,024}.
\end{align}
Counting another $VD$ for the output head would be incorrect because its
weights are tied to the input table.

At batch size 32 and context 128, each update contains
$32\cdot128=4096$ prediction targets before any proposed masking:
\begin{align}
 D_{\rm ppt}&=100\cdot4096=409\,600,\\
 D_{\rm language}&=900\cdot4096=3\,686\,400,\\
 D_{\rm total}&=1000\cdot4096=4\,096\,000.
\end{align}
Five arms and three seeds give
$15\cdot4\,096\,000=61\,440\,000$ total training tokens.
Those are historical design counts, not evidence that they are
matched-language-stage budgets.

The commonly used dense-parameter training proxy gives
\begin{equation}
 6P_{\rm model}D_{\rm total}
 =45\,638\,221\,824\,000
 \approx4.56382\times10^{13}
\end{equation}
arithmetic operations per run. The factor six approximates two
operations per weight for forward multiplication and four for backward
work. It is a proxy, not an exact FLOP count: attention's token-pair
work, generator operations, nonlinearities, memory movement and runtime
overhead are not all represented. Equal parameter/token products do not
establish equal measured time.

\subsection{Interpreting recorded losses mathematically}

For loss $L$ in nats per byte, byte perplexity is $\exp(L)$.
For paired losses, the geometric perplexity ratio follows from a mean
loss difference:
\begin{equation}
 \exp\!\left(\frac1S\sum_{s=1}^{S}(L_{G,s}-L_{0,s})\right).
\end{equation}
The recorded pilot mean test difference of approximately $0.053230$
nats/byte gives $\exp(0.053230)\approx1.05467$, about $5.47\%$ worse
geometric-mean byte perplexity. It is not a $5.47\%$ accuracy difference.

For positive finite paired cost ratios $r_s=C_{G,s}/C_{0,s}$, the
paired geometric ratio is
\begin{equation}
 R=\exp\!\left(\frac1S\sum_s\log r_s\right).
\end{equation}
The recorded historical ratio $R\approx1.432$ corresponds to a speedup
$R^{-1}\approx0.6983$, favoring scratch for that declared pipeline metric.
It neither resolves the scheduler confound nor identifies a mechanism.
An arithmetic ratio of group mean times generally differs from $R$.

If a target is not reached, its first-hit cost is censored. Removing
those seeds and averaging only finite ratios changes the population
being summarized. No finite universal speedup follows from that
operation. A fixed evaluation cadence also locates the first true
crossing only within an interval if loss is monotone between evaluations;
without monotonicity, an earlier unobserved crossing cannot be excluded.
The toy quadratic is monotone; real validation curves need not be.

